\subsection{An actual lift through invariant blocks with free cusp characters}

\begin{statement}{D.59}{Proved}{Further arithmetic invariant central coordinates lift by an integral homology correction and an injective norm}
Retain the actual arithmetic source and all maps of D.58.
Write
\[
 \Delta=N/K_3\simeq(\mathbb Z/3)^{97},\qquad
 G/\Delta=Q,\quad |Q|=576,\qquad
 \mathcal K=\ker(Z\longrightarrow Z_*).
 \tag{D.59.1}
\]
The normal subgroup \(A\) of D.49 lies in \(\Delta\).
The 48 oriented cusp classes of the actual marked
rank-97 group \(N\) are denoted by \(c_j\in\Delta\).
They span a 47-dimensional subspace \(\mathcal R_N\).
Define the affine character space used in the certificate
and the full character family used in the lift by
\[
 \begin{aligned}
 \mathscr A_c&=\{\ell\in\operatorname{Hom}(\Delta,\mathbb F_3):
                          \ell(c_j)=1\ (1\leq j\leq48)\},\\
 \Gamma&=\{\ell:\ \ell(c_j)\ne0\text{ for every }j,
          \ \ell|_A\ne0,\\
          &\hspace{19mm}\operatorname{Stab}_Q\{\ell,-\ell\}=1\},\\
 \kappa&=|\{\{\ell,-\ell\}:\ell\in\Gamma\}/Q|.
 \end{aligned}
 \tag{D.59.2}
\]
The affine space has dimension 50. Negation sends
its cusp values to two, so no \(Q\)-orbit in
\(\mathscr A_c\) meets its negative. The actual
calculation of the 575 nonidentity fixed affine
systems gives
\[
 \begin{gathered}
 B_c:=\sum_{1\ne q\in Q}|\mathscr A_c^q|
                  =11\,351\,380\,622\,391,\\
 \kappa\geq
 \left\lceil{2\cdot3^{49}-B_c\over576}\right\rceil
          =830\,900\,448\,697\,714\,719\,719>0.
 \end{gathered}
 \tag{D.59.3}
\]
This bound uses the actual nonzero compact image
of \(f_h\) established in D.56; it does not
evaluate the entire arithmetic subgroup \(A\).

Here are the projectors and the source they produce.
Put \(\mathcal O=\mathbb Z_2[\zeta_3]\), and, for
\(\ell\in\operatorname{Hom}(\Delta,\mathbb F_3)\), set
\[
 e_\ell={1\over|\Delta|}
             \sum_{d\in\Delta}\zeta_3^{-\ell(d)}d,
 \qquad e=\sum_{\ell\in\Gamma}e_\ell,
 \qquad \mathscr E=eU.
 \tag{D.59.4}
\]
The sum \(e\) belongs to \(\mathbb Z_2[\Delta]\)
and is central in \(\mathbb Z_2[G]\).
The integral module \(\mathscr E\) is projective
over \(\mathbb Z_2[G]\), has rank
\(110592\kappa\), and is contained in the exact
kernel \(\mathcal L=\ker(U^-\to U_*^-)\) of D.58.
The operator \(1-t\) is invertible on \(\mathscr E\).

On the central module define the actual equivariant
idempotent and its image by
\[
 \Pi=e_\Delta\bigl(\bigwedge\nolimits^2\bar e\bigr),
 \qquad \mathcal C=\Pi Z,\qquad
 \mathcal K_\Delta=\mathcal K^\Delta,\qquad
 \mathcal J=(1-\Pi)\mathcal K_\Delta,
 \tag{D.59.5}
\]
where \(\bar e\) is reduction modulo two and
\(e_\Delta\) is odd averaging on the exterior
module. These formulas give actual normal
subgroups of the full D.54 source
\(X=\widehat B/(R_{50}\cap L_2)\). In particular
\[
 \mathcal C\subset\mathcal K_\Delta,
 \qquad
 \mathcal C\simeq
       \mathbb F_2[Q]^{\,9216\kappa},\qquad
 \mathcal K_\Delta=\mathcal C\oplus\mathcal J.
 \tag{D.59.6}
\]

There is an actual continuous peripheral automorphism
of
\[
                         X_{59}=X/\mathcal J.
 \tag{D.59.7}
\]
It has multiplier one, fixes \(t\), induces identity
on \(G\), and retains the complete D.57 action
and its compatible conjugators on \(X_*\).
It retains the entire prescribed \(\mathcal W_h\)
action and its literal normalized parameter.
The full integral homology action is changed from
the chosen D.58 matrix \(T_0\) to the exact matrix
\[
                       T_1=(1-e)T_0+e.
 \tag{D.59.8}
\]
Thus it is identity on \(\mathscr E\), while its
action on the complementary integral summand is
unchanged. The already chosen D.58 homology
representative is replaced in this explicitly
specified summand.

The comparison to the same underlying source
\(Y=X/\mathcal K\) has exact kernel
\[
 \ker(X_{59}\to Y)
       =(1-e_\Delta)\mathcal K\oplus\mathcal C.
 \tag{D.59.9}
\]
The action induced on \(Y\) is the repaired action
with matrix (D.59.8). The remaining full class-two
map is \(X\to X_{59}\), with exact kernel
\(\mathcal J\). Its dimension is
\[
 \dim\mathcal J=48w-2560-5\,308\,416\kappa,
 \qquad w=96\cdot3^{97}(1-|A|^{-1}).
 \tag{D.59.10}
\]
It contains a specified 366080-dimensional submodule
coming from the nonzero characters pulled back
from \(A_*\). This residual is therefore nonzero.
The full source kernel is precisely
\[
 S_{59}=\{b\in R_{50}:
               \operatorname{im}_{\mathcal T_2}(b)\in\mathcal J\},
 \qquad S_{59}/(R_{50}\cap L_2)=\mathcal J.
 \tag{D.59.11}
\]
The remaining invariant boundary in
\(\mathcal J_G\) and the full profinite realization
are not evaluated by this statement.
The manuscript passages are Part D(e), section
51.16, folio 734 verso, scan 174, and section
51.15, folios 725--727, scans 156--160.
\end{statement}
\begin{proof}
\emph{The actual cusp-character calculation.}
The marked \(Q\)-cover has cusp width 12 and
48 cusps. Its free kernel \(N\) has rank 97
and compact genus 25. The positive cusp words
are the actual conjugates of \(t^{12}\).
Their homology classes have their one sum
relation and span a primitive rank-47 integral
submodule. Reduction modulo three gives the
displayed \(\mathcal R_N\), and assigning the
value one to all of them is consistent because
\(48=0\) in \(\mathbb F_3\). Thus
\(\dim\mathscr A_c=97-47=50\). The action of
\(Q\) permutes these oriented classes, proving
the asserted invariance of the affine space
and its separation from its negative.

The program \texttt{d\_arithmetic\_free\_characters.py}
reconstructs the actual 576-state graph and
all 576 matrices on \(N^{\rm ab}/3\). It checks
all 1152 generator matrix products. It forms
the 48 actual vectors \(q[t^{12}]\), verifies
their rank 47, their sum relation and both
complete generator permutations. It then
solves all 575 affine fixed systems over
\(\mathbb F_3\). The largest consistent fixed
dimension is 26, and the sum of their exact
cardinalities is \(B_c\) in (D.59.3). The
certificate retains every system's consistency
and dimension, the full affine parametrization,
and 51 explicitly recorded free characters
whose affine span has dimension 50.

Let \(a=f_h\bmod G\), viewed in \(\Delta\).
D.56's actual quotient
\(\phi:\Delta\twoheadrightarrow A_*\) kills
every one of these cusp classes: its cusp
still has order 12. Its value on \(a\) is
the nonzero vector \((2,1,2,1)\).
The program checks these 48 zero images
against the frozen actual quotient matrix.
Consequently \(a\notin\mathcal R_N\).
Evaluation on \(a\) is therefore a nonconstant
affine function on \(\mathscr A_c\); precisely
\(2\cdot3^{49}\) of its characters have
\(\ell(a)\ne0\). They are active on \(A\).
The union of all nonidentity fixed loci has
size at most \(B_c\). Every free active orbit
in \(\mathscr A_c\) has 576 elements and
gives a distinct unordered-pair orbit in
\(\Gamma\), since its negative has cusp
values two. These exact maps and the
union bound prove (D.59.3). No unknown
coordinates of \(a\) are assigned in this
argument.

\emph{Integral projectors and projectivity.}
The Fourier idempotents in (D.59.4) are defined
over the unramified quadratic ring \(\mathcal O\),
because \(|\Delta|\) is odd. Frobenius sends
\(e_\ell\) to \(e_{-\ell}\), so their indicated
sum is integral over \(\mathbb Z_2\).
The set is \(Q\)-stable, proving centrality.
Each included character is nontrivial on
\(A\), whence \(e_A\mathscr E=0\).

Here is the exact multiplicity calculation.
The abelian three-cover of the free rank-97
group \(N\) has its usual marked graph sequence
over \(\mathbb Z_2[\Delta]\):
\[
 0\longrightarrow U\longrightarrow
       \mathbb Z_2[\Delta]^{97}
       \longrightarrow I_\Delta\longrightarrow0.
 \tag{D.59.12}
\]
On a nontrivial character the last map is a
surjection \(\mathcal O^{97}\to\mathcal O\):
at least one of its coefficients is a unit
\(\zeta_3^j-1\). Its kernel is free of rank
96 over \(\mathcal O\). Fix an unordered
pair \(\{\ell,-\ell\}\) from (D.59.2).
Its integral isotypic module has rank 192
over \(\mathbb Z_2\). Since its stabilizer
in \(Q\) is trivial, the literal map
\[
 \operatorname{Ind}_\Delta^G
          U_{\{\ell,-\ell\}}\longrightarrow
             U_{Q\{\ell,-\ell\}},
 \qquad g\otimes v\longmapsto gv
 \tag{D.59.13}
\]
is an isomorphism: its summands are the
distinct isotypic components. The source
module over \(\Delta\) is a direct summand
of copies of \(\mathbb Z_2[\Delta]\), so
its induced module is projective over
\(\mathbb Z_2[G]\). Its rank is
\(576\cdot192=110592\), proving all
projectivity and rank assertions for
\(\mathscr E\).

The element \(t^{12}\in\Delta\) acts on
every character component of \(\mathscr E\)
as either \(\zeta_3\) or \(\zeta_3^2\).
Thus \(1-t^{12}\) is invertible. The
identity
\[
 (1-t)(1+t+\cdots+t^{11})=1-t^{12}
 \tag{D.59.14}
\]
shows that \(1-t\) is invertible too.
In the actual small quotient \(G_*\), the
element \(t^{12}\) is one. It therefore
acts as identity on \(U_*\). Equivariance
of \(\theta_U\) and (D.59.14) give
\(\theta_U\mathscr E=0\), with the
actual source-to-target map. Together with
\(e_A\mathscr E=0\), this proves
\(\mathscr E\subset\mathcal L\).

\emph{Every prescribed homology correction in this summand is peripheral.}
For these integral coefficients the same
averaging argument as D.53 gives
\(H^i(G,\mathscr E)=0\) for every \(i\geq0\):
the trivial module is in the \(e_A\) block,
while \(\mathscr E\) is in its complementary
block. Applying the central idempotent to a
free resolution proves the assertion over
\(\mathbb Z_2\) as well as modulo two.
Inflation--restriction therefore gives the
actual restriction isomorphism
\[
 H^1(\widehat B,\mathscr E)
       \xrightarrow{\ \sim\ }
          \operatorname{Hom}_G(U,\mathscr E).
 \tag{D.59.15}
\]
The restriction to \(K_3\) factors through
its maximal abelian pro-two quotient \(U\).
The induced module in (D.59.13) restricts
projectively to \(\langle s\rangle\), so
\(H^1(\langle s\rangle,\mathscr E)=0\).
For \(\langle r\rangle\), odd averaging
gives the same vanishing. Hence a derivation
representing any prescribed map on the
right of (D.59.15) is an actual coboundary
on each finite peripheral subgroup.

Its cusp value is removed by subtracting
the inner derivation \((1-g)b\), where
\(b=(1-t)^{-1}\delta(t)\). This does not
change its restriction to \(K_3\). It
follows that every map in (D.59.15) is
represented by an integral derivation
with
\[
 \delta(s)=(1-s)b_s,\qquad
 \delta(r)=(1-r)b_r,\qquad \delta(t)=0.
 \tag{D.59.16}
\]
These equations define actual continuous
homomorphisms into the compact semidirect
product; they are not merely infinitesimal
corrections.

Apply this construction to
\((1-T_0)e:U\to\mathscr E\). The kernel
\(Y\to X_*\) is the actual integral module
\(\mathcal L\) by D.58. Thus (D.59.16)
changes the two conjugators in that kernel
and induces the exact homology matrix
(D.59.8). Its two integral blocks are
invertible. The actual maps to \(X_*\)
and to the full homology quotient agree
on their common quotient, since the
corrections lie in \(\mathcal L\).
Their fibre product therefore gives an
actual automorphism of \(Y\), retaining
the fixed D.57 conjugators. This proves
the required integral correction before
any central equation is used.

\emph{The invariant exterior module and its compact image.}
The operator \(\Pi\) in (D.59.5) is an
idempotent, as \(\bar e\) and odd averaging
commute with the \(G\)-action. Over
\(\mathbb F_4\), a \(\Delta\)-invariant
exterior term can only pair a character
with its negative. Since every character
under consideration is active, these
terms belong to \(Z=W\wedge V\).
The map to \(Z_*\) kills them because
\(\theta_V\bar e=0\). This proves
\(\mathcal C\subset\mathcal K_\Delta\)
and the direct-sum decomposition in
(D.59.6) by the displayed projector.

For a single unordered character pair,
choose \(96\) pairs of \(\mathbb F_2\)
vectors \(u_i,v_i\), on which a nontrivial
element of the character group acts by
\((u_i,v_i)\mapsto(v_i,u_i+v_i)\).
Its invariant exterior module has the
explicit basis
\[
 \begin{gathered}
 u_i\wedge v_i,\\
 u_i\wedge u_j+u_i\wedge v_j+v_i\wedge v_j,
 \qquad u_i\wedge v_j+v_i\wedge u_j\quad(i<j).
 \end{gathered}
 \tag{D.59.17}
\]
It has dimension \(96^2=9216\).
There are 576 distinct pairs in the
orbit, with no stabilizer. The literal
translation map from a chosen fibre
therefore identifies its exterior module
with \(\mathbb F_2[Q]^{9216}\).
This proves (D.59.6), including its
actual action through \(G\to Q\).
The full group norm
\[
 \nu:\mathcal C_G\longrightarrow\mathcal C^G,
       \qquad[z]\longmapsto\sum_{g\in G}gz
 \tag{D.59.18}
\]
is an isomorphism: \(\Delta\) acts
trivially and has odd order, while on
a regular \(Q\)-module the map sends
one fibre to its diagonal sum of all
576 translates.

The compactification map of D.55 is
injective on the entire homology
summand \(\bar eV\). Indeed its cusp
module is a quotient of
\(\mathbb F_2[G/\langle t\rangle]\).
The intersection
\(\Delta\cap\langle t\rangle\)
is the order-three group generated by
\(t^{12}\). Every character in
\(\Gamma\) is nontrivial
on every conjugate of that group.
Its projector consequently kills the
entire cusp module. Thus the actual
compactification map
\(\pi:V\to H\) identifies \(\bar eV\)
with its image, and
\(\bigwedge^2\pi\) is injective on
\(\mathcal C\). This removes the
compact-projection kernel in (D.55.5)
for exactly the blocks used here.

\emph{The central boundary vanishes in the specified new coordinates.}
Lift the repaired \(Y\) conjugators
to \(X\). Their cusp error is an
actual element \(z\in\mathcal K\).
The induced homology map is \(T_1\).
Apply the exact orientation identity
of D.55 and then the projector induced
by \(\Pi\) on compact exterior homology.
The right side is zero, because
\(T_1\) is identity on \(\mathscr E\).
The left side is
\((\bigwedge^2\pi)\nu([\Pi z])\).
Both maps are injective on this
specified block by (D.59.18) and the
preceding paragraph. Therefore
\[
                          [\Pi z]=0\quad\hbox{in }\mathcal C_G.
 \tag{D.59.19}
\]
The other retained central summand
\((1-e_\Delta)\mathcal K\) has zero
\(\Delta\)-coinvariants. Its cusp
error is consequently in its
augmentation image as well. In both
summands that image is exactly
\((1-s)M+s(1-r)M\). Actual central
corrections to the two conjugators
kill these two errors and preserve
their projection to \(Y\). The
remaining error lies in \(\mathcal J\),
so the cusp identity is exact in
\(X/\mathcal J\).

The matrix \(T_1\) commutes with
all displayed projectors and preserves
\(\mathcal K\), since its small
projection is unchanged. Its exterior
square therefore preserves
\(\mathcal J\). The corrected source
map descends to the actual finite
quotient \(\mathcal T_2/\mathcal J\)
and is an automorphism by its invertible
Frattini map. It agrees with the
repaired automorphism of \(\mathcal Z_h\),
induced by its projections to \(\mathcal W_h\)
and the full homology quotient, on
their original common quotient.
The fibre product is exactly
\(X/\mathcal J\), proving the claimed
continuous automorphism and the exact
kernel (D.59.9). All changes are in
the source kernel over \(X_*\);
the parameter and CM conclusions
follow through the unchanged
\(\mathcal W_h\) projection.

\emph{Dimensions and an explicit remaining central submodule.}
Reducing (D.59.12) modulo two and
using semisimplicity for the odd
group \(\Delta\) gives
\[
 V|_\Delta\simeq
           \mathbb F_2\oplus\mathbb F_2[\Delta]^{96}.
 \tag{D.59.20}
\]
There are
\(|\Delta|(1-|A|^{-1})\) active
nontrivial characters over
\(\mathbb F_4\). Pairing each with
its negative, with the multiplicity
96 just proved, gives
\[
 \dim Z^\Delta
       ={96^2\over2}|\Delta|(1-|A|^{-1})=48w.
 \tag{D.59.21}
\]
Odd averaging is exact and the
target action factors through
\(\Delta\twoheadrightarrow A_*\).
Thus \(Z^\Delta\to Z_*^{A_*}\)
is onto, with target dimension 2560.
Subtracting that dimension and the
dimension \(576\cdot9216\kappa\)
of \(\mathcal C\) proves (D.59.10).

For an explicit remaining submodule,
take all 40 nonzero unordered
characters pulled back from \(A_*\).
They are active, since \(A\to A_*\)
is onto. Their values on every
\(c_j\) are zero, so their central
projector is disjoint from \(\Pi\).
The source character multiplicity
is 96; the target multiplicity is
eight, by the rank-nine cover in
D.57. The actual exterior map is
onto by odd averaging and the
onto homology map. Its kernel in
these specified components has
dimension
\[
                   40(96^2-8^2)=366080.
 \tag{D.59.22}
\]
It is a \(G\)-submodule of
\(\mathcal J\). For a direct
nonzero coordinate, choose a
nonzero vector \(u\) in the
underlying \(\mathbb F_2\)-kernel
of one unordered-pair homology map, and
choose \(d\in\Delta\) with
character value one. On its
underlying two-dimensional
\(\mathbb F_2\)-line, the element
\(u\wedge du\) is nonzero and
\(\Delta\)-invariant. It maps
to zero in the small quotient
and belongs to the indicated
residual submodule. The actual
onto map \(R_{50}\to Z\) supplies
a source preimage. Its image in
\(X\) is nonzero and its image
in \(X_{59}\) is one. This proves
the residual assertion through
the literal quotient maps in
(D.59.11), without asserting a
full profinite lift.
\end{proof}
